\documentclass[12pt,a4paper]{article}
\usepackage[utf8]{inputenc}
\usepackage[spanish]{babel}
\usepackage{geometry}
\usepackage{listings}
\usepackage{xcolor}
\usepackage{parskip}
\usepackage{amsmath}
\usepackage{hyperref}

\geometry{margin=2.5cm}

\lstset{
  basicstyle=\ttfamily\small,
  keywordstyle=\color{blue}\bfseries,
  commentstyle=\color{gray},
  breaklines=true, frame=single, language=Java,
  numbers=left, numberstyle=\tiny\color{gray}
}

\title{\textbf{Problema de Rutas --- BFS, UCS y A* (v2)}\\
  \large Documentación Técnica}
\author{Inteligencia Artificial}
\date{}

\begin{document}
\maketitle
\tableofcontents
\newpage

\section{Descripción del Proyecto}
Versión extendida del planificador de rutas que agrega el algoritmo \textbf{A*} con soporte de heurísticas. Implementa BFS, UCS y A* sobre el mismo grafo de ciudades mexicanas.

\section{Algoritmo A* (A-Star)}

\subsection{Fundamento}
A* combina el costo real acumulado $g(n)$ con una estimación heurística $h(n)$ de la distancia restante al destino:
\[
f(n) = g(n) + h(n)
\]

Para que A* sea óptimo, la heurística debe ser \textbf{admisible}: nunca sobreestimar la distancia real.

\subsection{Heurística utilizada}
Se usa la distancia en línea recta (distancia euclidiana) entre ciudades, calculada con sus coordenadas geográficas aproximadas:

\begin{lstlisting}
double heuristica(String ciudad, String destino) {
    double[] coordCiudad  = coordenadas.get(ciudad);
    double[] coordDestino = coordenadas.get(destino);
    double dLat = coordCiudad[0] - coordDestino[0];
    double dLon = coordCiudad[1] - coordDestino[1];
    // factor 111 km por grado de latitud
    return Math.sqrt(dLat*dLat + dLon*dLon) * 111;
}
\end{lstlisting}

Esta heurística es admisible porque la línea recta es siempre menor o igual que la ruta real por carretera.

\subsection{Implementación A*}
\begin{lstlisting}
Nodo resolverAEstrella(String origen, String destino) {
    PriorityQueue<Nodo> abiertos =
        new PriorityQueue<>(Comparator.comparingDouble(
            n -> n.g + heuristica(n.ciudad, destino)));
    Set<String> cerrados = new HashSet<>();

    abiertos.add(new Nodo(origen, 0, null));

    while (!abiertos.isEmpty()) {
        Nodo actual = abiertos.poll();
        if (cerrados.contains(actual.ciudad)) continue;
        cerrados.add(actual.ciudad);

        if (actual.ciudad.equals(destino)) return actual;

        for (Arista a : grafo.get(actual.ciudad)) {
            if (!cerrados.contains(a.destino)) {
                double nuevoG = actual.g + a.peso;
                abiertos.add(new Nodo(a.destino,
                                      nuevoG, actual));
            }
        }
    }
    return null;
}
\end{lstlisting}

\section{Comparativa de los Tres Algoritmos}

\begin{center}
\begin{tabular}{|l|c|c|c|}
\hline
\textbf{Característica} & \textbf{BFS} & \textbf{UCS} & \textbf{A*} \\
\hline
Estructura & Cola FIFO & Cola prioridad & Cola prioridad \\
Criterio & Saltos & $g(n)$ & $f(n)=g(n)+h(n)$ \\
Heurística & No & No & Sí \\
Óptimo en costo & No & Sí & Sí \\
Nodos explorados & Más & Medio & Menos \\
Velocidad & Media & Lenta & Rápida \\
\hline
\end{tabular}
\end{center}

\section{Clase ResultadoBusqueda}

Para unificar la salida de los tres algoritmos se usa una clase resultado:

\begin{lstlisting}
class ResultadoBusqueda {
    List<String> ruta;
    double costoTotal;
    int nodosExplorados;
    long tiempoMs;
}
\end{lstlisting}

Esto permite mostrar comparativas de rendimiento entre los tres algoritmos para la misma consulta.

\section{Interfaz Gráfica}

\begin{itemize}
  \item Tres botones: \textbf{BFS}, \textbf{UCS}, \textbf{A*}.
  \item Muestra ruta, costo y nodos explorados para comparar eficiencia.
  \item La ruta se visualiza resaltada sobre el grafo de ciudades.
\end{itemize}

\end{document}
